\documentclass[10pt,a4paper]{amsart}
\usepackage[margin=19mm]{geometry}
\usepackage{amsmath,amssymb,mathtools}
\usepackage[scr=rsfs,cal=euler]{mathalfa}
\usepackage{xfrac}
\usepackage[unicode,hidelinks,bookmarks=false]{hyperref}
\newcommand{\ie}{i.e.\ }
\newcommand{\cf}{cf.\ }
\newcommand{\resp}{respectively\ }
\newcommand{\Subsection}{\subsection}
\providecommand{\nobreakdash}{\nobreak\hbox{-}}
\setlength{\emergencystretch}{2em}
\multlinegap0pt
\usepackage{tikz}
\usetikzlibrary{cd}
\usepackage{amssymb}
\usepackage{xcolor}
\usepackage{tikz}
\newtheorem{theorem}{Theorem}
\newtheorem{remark}{Remark}
\newtheorem{proposition}{Proposition}
\newtheorem{lemma}{Lemma}
\newcommand{\DM}[2]{\mathbf{DM}_{#2}^{\mathit{eff}}(#1)}
\newcommand{\Z}{{\mathbb Z}}
\newcommand{\hocolim}{{\rm hocolim \,}}
\newcommand{\sH}{{\mathcal H}}
\newcommand{\sW}{{\mathcal W}}
\newcommand{\sX}{{\mathcal X}}
\newcommand{\sY}{{\mathcal Y}}
\newcommand{\sZ}{{\mathcal Z}}
\title{On the motivic homotopy type of algebraic stacks}
\author{Neeraj Deshmukh, Jack Hall}
\date{Source: arXiv:2304.10631v4 (CC BY 4.0)}
\begin{document}
\maketitle

\noindent\emph{Source and licence.} On the motivic homotopy type of algebraic stacks. Version arXiv:2304.10631v4 (CC BY 4.0), distributed by arXiv under the Creative Commons Attribution 4.0 International licence, http://creativecommons.org/licenses/by/4.0/, as recorded in the arXiv licence metadata for this exact version and shown on the abstract page https://arxiv.org/abs/2304.10631v4. Full licence notice: \texttt{licenses/arxiv-2304.10631-CC-BY-4.0.txt}.\par\smallskip

\noindent\textit{Editorial scope.} Counted unit: the article's proposition nisnevich-descent (source0.tex lines 1466--1474). The article's complete original proof of it is delivered with it (source0.tex lines 1476--1514, one \texttt{proof} environment of 39 lines). No mathematics is written, repaired or re-derived: the statement and the proof are the article's own lines, in the article's own notation, and no retained line is substituted or re-flowed.  Retained uncounted as the source-local context the proof needs: lines 1197--1199; lines 1272--1275; lines 1453--1464; lines 1739--1743.  Nothing was omitted from any retained span, so the package carries no omission record.  No retained line contains a citation, so the wrapper carries no bibliography and no bibliography entry.  The retained lines contain the article's own diagram code, which the wrapper typesets with the standard tikz packages; no image file is delivered and the package declares no asset dependency.  Editorial formatting only, in addition: an AMS \texttt{amsart} preamble that restates the article's own declarations and macros used by the retained lines, namely the environments \texttt{theorem}, \texttt{remark}, \texttt{proposition}, \texttt{lemma} on separate counters, together with the standard packages those lines need.  The retained environments print in this document as Theorem 1, Remark 1, Proposition 1, Lemma 1 in order of appearance, whereas the article prints the same items with the numbering of its own full document.  The complete native source of the article as distributed in the arXiv e-print for this exact version is bundled unchanged as \texttt{sources/140378/source0.tex}.
\begin{theorem}\label{motive-construction}
	Let $\sY\rightarrow \sX$ be a representable Nisnevich cover of algebraic stacks over a field $k$. Assume that $\sY$ is of the form $[Y/GL_n]$ for some algebraic space $Y$, and let $p:Y\rightarrow\sY\rightarrow\sX$ be the composite. Then, for the \v{C}ech hypercover $Y_{\bullet}$ associated to $p $, $p_{\bullet}:Y_{\bullet}\rightarrow \sX$ is a Nisnevich local weak equivalence, i.e, $Y_{\bullet}\simeq \sX$ in $\sH_{\bullet}(k)$.
\end{theorem}

\begin{remark}
	\label{remark-motive-commutes-with-homotopy-colimits}
	$N\Z_{tr}$ is a left Quillen functor with respect to the Nisnevich motivic model structures on $\mathcal{H}_{\bullet}(k)$ and $\DM{k,\Z}{}$. Thus, Lemma \ref{homotopy-colimits-left-adjoints} implies that $M$ commutes with homotopy colimits.
\end{remark}

\begin{remark}\label{BaseChange}
Let $\sZ$ be a cd-quotient stack. If $\sY\rightarrow\sZ$ is a representable morhpism, then $\sY$ is also a cd-quotient stack. To see this, let $Z\rightarrow [Z/GL_n]\rightarrow\sZ$ be a Nisnevich covering of $\sZ$ by a quotient stack. Now, observe that we have a cartesian diagram by base change,
\begin{center}
\begin{tikzcd}
	\sY\times_{\sZ}Z\arrow[r]\arrow[d] & Z\arrow[d]\\
	\sY\times_{\sZ}[Z/GL_n]\arrow[r]\arrow[d] & \left[ Z/GL_n\right]\arrow[d]\\
	\sY\arrow[r] & \sZ
\end{tikzcd}
\end{center}
where $\sY\times_{\sZ}Z\rightarrow \sY\times_{\sZ} [Z/GL_n]$ is a $GL_n$-torsor and $\sY\times_{\sZ} [Z/GL_n]\rightarrow\sY$ is a Nisnevich cover. Denote the algebraic space $\sY\times_{\sZ}Z$ by $Y$. Thus, $\sY$ is a cd-quotient stack. Hence, by Theorem \ref{motive-construction}, $Y_{\bullet}\rightarrow \sY$ is a Nisnevich local weak equivalence.\\
This tells us that \textit{$GL_n$-presentations} respect base change.
\end{remark}

\begin{lemma}\label{homotopy-colimits-left-adjoints}
	Let $F:M\rightleftarrows N: G$ be a Quillen adjunction of model categories. Let $I$ be an index category such that the projective model structure is defined on $M^I$ and $N^I$. Then, for an $I$-diagram $E$ in $M$, we have
	\[LF(\hocolim E)\simeq \hocolim LF^I(E),\]
	where $LF$ and $LF^I$ are the derived functors of $F$ and $F^I$, repectively.
\end{lemma}

\begin{proposition}\label{nisnevich-descent} For a distinguished Nisnevich square,
\begin{center}
\begin{tikzcd}
	\sW\arrow[r]\arrow[d] & \sY\arrow[d,"p"]\\
	\sX\arrow[r,hook,"j"]	&\sZ
\end{tikzcd}
\end{center}
where $j$ is an open immersion and $p$ is \'{e}tale representable, the induced diagram on motives is homotopy cartesian.
\end{proposition}

\begin{proof}
For a distinguished Nisnevich square of stacks
\begin{center}
\begin{tikzcd}
	\sW\arrow[r]\arrow[d] & \sY\arrow[d,"p"]\\
	\sX\arrow[r,hook,"j"]	&\sZ
\end{tikzcd}
\end{center}
we need to show that the induced diagram 
\begin{center}
	\begin{tikzcd}
		M(\sZ)\arrow[r,,"p^*"]\arrow[d,"j^*"] & M(\sY)\arrow[d]\\
		M(\sX)\arrow[r,]	&M(\sW)
	\end{tikzcd}
\end{center}
of motives is homotopy cartesian.
If $Z\rightarrow [Z/GL_n]\rightarrow\sZ$ is a Nisnevich covering of $\sZ$ by a $GL_n$-torsor, then by Remark \ref{BaseChange}, we can base change this covering to $\sY,\sX$ and $\sW$. This gives us a cartesian diagram of hypercovers:
\begin{center}
\begin{tikzcd}
	W_{\bullet}\arrow[r]\arrow[d] & Y_{\bullet}\arrow[d,"p_{\bullet}"]\\
	X_{\bullet}\arrow[r,hook,"j_{\bullet}"]	&Z_{\bullet}
\end{tikzcd}
\end{center}
For each $i$, the above the diagram is a distinguished Nisnevich square of schemes. Thus, for each $i$, the following diagram of motives is homotopy (co)cartesian,
\begin{center}
\begin{tikzcd}
	M(Z_i)\arrow[r]\arrow[d] & M(Y_i)\arrow[d]\\
	M(X_i)\arrow[r]	&M(W_i)
\end{tikzcd}
\end{center}
By Remark \ref{remark-motive-commutes-with-homotopy-colimits}, $M(Z_{\bullet})\simeq \hocolim M(Z_i)$. As homotopy colimits commute with homotopy colimits, the following diagram is again homotopy (co)cartesian,
\begin{center}
\begin{tikzcd}
	M(Z_{\bullet})\arrow[r]\arrow[d] & M(Y_{\bullet})\arrow[d]\\
	M(X_{\bullet})\arrow[r]	&M(W_{\bullet})
\end{tikzcd}
\end{center}
By Theorem \ref{motive-construction}, we get the required result.
\end{proof}

\end{document}
